\documentclass[5p,times]{elsarticle}

\usepackage{booktabs}

\journal{---}

\begin{document}

%\begin{frontmatter}

\title{An Agentic Multi-Agent ML Pipeline for Predictive Maintenance with a Deterministic Trust Boundary}

\author[rit]{Michael Eaton\corref{cor1}}
%\ead{me3870@rit.edu}
%\address[rit]{Data Science, Rochester Institute of Technology, Rochester, NY, USA}

%\begin{abstract}
%Agentic AI systems built on Large Language Models (LLMs) offer contextual reasoning and adaptive tool use, but industrial machine learning (ML) applications such as predictive maintenance demand results that are verifiably correct, not merely plausible. This work presents a hybrid multi-agent classification pipeline in which an orchestrator coordinates six specialized LLM agents,intake, feature engineering, profiling, modeling, verification, and post-hoc deep-dive explanation,while a fully deterministic, non-LLM harness retains exclusive authority over data splitting, metric computation, leakage detection, and sandboxed execution. Agents make judgment calls as structured proposals; the harness independently re-validates every proposal before acting on it, so a hallucinated column, a leaky preprocessing choice, or an unearned claim of performance is rejected rather than propagated. The system is validated end-to-end on an aviation predictive maintenance dataset (NGAFID-MC), including bounded-memory featurization of multi-gigabyte sensor streams, entity-aware splitting by aircraft, and an evidence-grounded deep-dive agent that explains individual flagged flights. Two additional tabular benchmarks demonstrate generalization across binary and multiclass problems. The current orchestrator executes a fixed agent sequence; extending it to dynamic, planner-driven orchestration over streaming data is identified as future work.
%\end{abstract}

%\end{frontmatter}

\section{Introduction}
Recent hybrid frameworks combine LLM-based orchestration with specialized agents for prescriptive maintenance workflows~\cite{farahani2026}. These systems demonstrate that agentic reasoning can plan and execute complete analytical pipelines, but they inherit a structural risk: the most common ways an ML result turns out to be wrong, an incorrect train/test split, preprocessing fit on held-out data, a feature that copies the target, are precisely the errors an LLM can introduce and then confidently narrate past. The central design question addressed here is therefore not ``can agents run an ML workflow?'' but ``what is the smallest system in which agents can make real modeling decisions without being able to fool themselves, or us, about whether those decisions were any good?''

\section{System Design}
The architecture divides authority along a hard line (Table~\ref{tab:authority}). An orchestrator invokes six specialized agents in sequence, each with a minimal decision surface, one or two read-only tools, and a strict JSON output contract. Every agent output is a \emph{proposal} that a deterministic harness independently re-validates before acting on it.

\begin{table}[h]
\caption{Division of authority between agents and harness.}
\label{tab:authority}
\centering
\small
\begin{tabular}{p{3.6cm}p{3.6cm}}
\toprule
\textbf{Agents propose} & \textbf{Harness decides} \\
\midrule
Target column and task framing & Train/val/test splits (group- and time-aware) \\
Feature drops and stateless derived features & Imputation, scaling, all fitted statistics \\
Model template and hyperparameters & Metric computation with bootstrap CIs \\
Plain-language narration and audit verdicts & Leakage gates, sandboxing, candidate promotion \\
\bottomrule
\end{tabular}
\end{table}

The six agents are: (i) an \emph{intake agent} that converts a dataset and optional natural-language goal into a formal problem specification; (ii) a \emph{feature engineering agent} restricted to a vetted catalog of stateless, row-wise operations, provably safe to apply before any split exists; (iii) a \emph{profiler agent} that narrates deterministic dataset facts and a rule-based split recommendation it cannot override; (iv) a \emph{modeling agent} that selects from six verified pipeline templates and fills in configuration, it never writes code; (v) a \emph{verification agent} with strictly one-directional power: it reviews only candidates that already passed every deterministic gate and can flag or reject, never approve or unblock; and (vi) a \emph{deep-dive agent} that synthesizes hedged explanations of individual flagged predictions from three deterministic measurements it cannot alter.

Safety is layered: schema, column, and configuration validation; static AST rejection of forbidden imports and calls; subprocess sandboxing with resource limits; a label-permutation test that catches leakage in the modeling \emph{process}; and a candidate-scoped correlation gate that catches leakage in raw feature \emph{content}. A dedicated test constructs a near-copy-of-target feature and demonstrates that each gate catches a failure the other misses, the two are complementary by construction, not redundant.

\section{Predictive Maintenance Use Case}
The pipeline was validated on the NGAFID-MC aviation dataset: 22 engine-sensor channels sampled at $\sim$1\,Hz across thousands of flights and multiple aircraft, with a binary before/after-maintenance target. A streaming featurizer rolls raw timesteps into one feature row per flight with memory bounded by chunk size rather than file size (4.2\,GB source), and entity-aware splitting by aircraft prevents the same airplane from appearing on both sides of the evaluation boundary. For flagged flights, the deep-dive agent combines flight-phase segmentation, model-agnostic occlusion attribution, and an independent raw-signal cross-cylinder localization check, then reports whether the model's attribution and the raw-signal evidence agree,a lead for inspection, not a diagnosis. Two tabular benchmarks (Titanic, Iris) additionally exercised binary and multiclass generalization; 87 automated tests, including full-loop integration tests with a stubbed model client, keep the entire agentic loop verifiable without network access.

\section{Limitations and Future Work}
The current orchestrator is a fixed sequence of phases: the same workflow executes on every run, and the deep-dive agent is invoked manually rather than routed to automatically. Future work targets the dynamic orchestration this architecture is designed to support: an LLM planner agent that selects which agents and workflows to instantiate per problem; ingestion agents for streaming sensor data arriving in real time; drift-monitoring and model-update agents that retrain and re-gate candidates as conditions change; cross-run evidence reuse; and parallel candidate search. Each extension inherits the same constraint that makes the present system trustworthy: new capability is added as a narrow, vetted agent behind the existing trust boundary, never by widening any agent's authority.

\section{Conclusion}
Hybrid agentic systems for manufacturing analytics are typically motivated by adaptability; this work argues that their adoption will equally depend on \emph{verifiability}. Placing every LLM judgment behind a deterministic harness that owns splitting, scoring, and leakage detection yields a system where agents contribute genuine modeling decisions while the correctness of the result never depends on an LLM being right.

\begin{thebibliography}{1}
\bibitem{farahani2026}
M.~A. Farahani, M.~I. Khan, T.~Wuest, Hybrid Agentic AI and Multi-Agent Systems in Smart Manufacturing, preprint submitted to the 2026 NAMRC Conference (2026).
\end{thebibliography}

\end{document}